\section{Comparison of the complete supplied corpus}
\label{cmp:comparison}

\subsection{Scope, identities and evidence levels}
\label{cmp:scope}

The comparison is pinned to ProveIt commit
\path|1c9524d8bf2e0a27b9a16ff622359927f536e663|.  It covers every delivered
research packet present under the requested PA research tree and
\path|docs/incoming| at that revision: 61 filed archives and 11 current ZIP
archives, hence 72 deliveries.  The accompanying ledger accounts for every
delivery separately.  The many adjacent formal-development directories
usually develop different components of the same machine theorem; their
directory count is not a count of different Turing machines.
The snapshot and its incoming reports are the sources of the historical
claims below.\cite{repo,incoming}

The fresh recognizer reads complete numeric binary transition tables,
including explicit JSON representations, and traverses retained nested
ZIPs.  It compares their actions up to control renaming and, separately,
global reflection.  In the current 194-state population there are
\emph{83 stored fingerprints}: 82 action graphs having a root from which
every state is reachable, plus one disconnected diagnostic fixture.
Reflection merges none of those 83 objects.  All 535 representatives in
the repository's earlier official catalogue were independently rechecked;
their fingerprints and recorded reachable-state counts all match.
The new machine constructed in this report is not included in these
counts of the supplied comparison corpus.

\begin{table}[htbp]
\centering
\begin{tabular}{@{}lr@{}}
\toprule
Historical evidence category at 194 states & Objects\\
\midrule
Reported principal construction & 40\\
Supplementary construction & 7\\
Research auxiliary, including source-tree rotations & 25\\
Research seed & 4\\
Explicitly rejected design & 6\\
Disconnected diagnostic & 1\\
\midrule
Total & 83\\
\bottomrule
\end{tabular}
\caption{An accounting of supplied tables, not a count of proved consistency
theorems.  Every object appears in Appendix~\ref{cmp:frontier-appendix}.}
\label{cmp:evidence-counts}
\end{table}

The distinction between those categories matters.  A packet may contain its
principal construction, several comparison inputs, experimental quotients,
and a synthetic test table.  A directory or filename containing ``193'' or
``194'' does not certify a consistency machine.  Likewise, a recorded
symbolic contract checks a particular interface; it does not establish a
logical calibration merely by being executable.  The historical categories
were matched to freshly parsed tables, but the old mathematical proofs and
all old certificate suites have not been independently rederived for this
report.  The new machine's checks are identified separately.

The scope is explicit rather than unlimited.  Recognized encodings include
five-column literal tables, textual tables embedded in JSON, transition
arrays, numeric-state dictionaries, transition-key dictionaries and row
records.  Arbitrary generators are not executed, and an unmaterialized
parameterized source program is not silently counted as an explicit table.
Source-only proposals remain represented in the delivery ledger.  Repeated
identical nested ZIP payloads are represented by duplicate references.
Files removed by the repository's earlier deduplication are accounted for
by its own catalogue and omission ledger, not reconstructed as fictitious
additional files in the fresh census.

A table identity also differs from a blank-start machine identity.  The
comparison may forget the start when testing graph isomorphism, but a
machine theorem must retain its designated start.  Some native tables
renumber their states from zero although zero is not their actual start.
For a graph with no spanning root, the census labels only its maximal
reachable components; this qualification is exposed in the data.  The
one disconnected 194-state diagnostic is not promoted to a complete
spanning graph class.

\subsection{What the principal constructions actually change}
\label{cmp:mechanisms}

The principal 194-state designs have extensive common ancestry: a
numerical proof calculus, finite command-word search, a unary register
file on a binary tape, and a shared source compiler.  Their substantive
differences occur in four places: the formulas emitted as seeds, the
MP/Shift command protocol, the binary decoder grammar, and context-sensitive
physical returns.  The earlier U208 and PA196 developments remain relevant
as sources of logical and implementation lemmas, but their larger state
counts do not themselves resolve the current target.\cite{repo,ax11packet}

\paragraph{Separate heads and the first 194-state machines.}
The original PA195 and its successors replace or reorganize the
diagonal-vacuity interface, then share the resulting axiom emitter with the
rest of the runtime.  Primary PA194 retains the separate diagonal head and
the original sum-based Shift test.  Its allocation is
\[
                  134+67-7=194:
\]
134 fixed controls, 67 raw axiom readers and seven simultaneous folds.
An alternate association uses $134+68-8=194$.  The two are different
action graphs with the same broad ten-entry AX interface.  An additional
A/D physical-register placement is another concrete variant.  These
changes must be distinguished from the many later files that merely
reproduce one of these tables in native or dense numbering.

\paragraph{Changing the command store.}
The C family changes successful modus ponens to place its result in the
cache and uses a cache-only Shift test.  CS adds a symmetric destructive
comparison that eliminates a transfer loop.  N instead omits the dummy-zero
push after successful MP, retains the conclusion visibly, and retains the
original sum guard.  Its proof compiler must reserve the required zero
explicitly; deleting the push while keeping the old proof words is not a
correctness argument.  H combines the visible no-push result with the
cache-only Shift test, and HS adds symmetric comparison.  The separate
SUMS and sum-guard NS machines use symmetric comparison while preserving
the original sum guard.  Thus ``no-push'', ``cached result'', ``cache-only
test'' and ``symmetric comparison'' are different choices, even when
several are composed in one machine.

These distinctions change what a refutation compiler must do.  A visible
conclusion one does not by itself trigger a cache-only Shift.  It must
first reach the cache, or the compiler must use a word whose internal
Shift tests it there.  Conversely, the original sum test does not halt
when both visible entries equal one.  Correct all-store proof words
therefore require exact command equations, including early-halt cases;
matching state counts or successful sample runs is inadequate.

\paragraph{Changing emitted formulas.}
Q exchanges the two curried premises in the universal-distribution head.
Finite K/S/MP proof translations justify the exchange; the change affects
a recurring block of emitted formulas, not only labels.  Q has been
combined with C, N, H and symmetric variants.  D instead replaces the
separate diagonal head $\Gamma(A)$ by $\Gamma(D)$, requiring a corresponding
selector and proof-word compiler.  DHS combines that emitter with no-push
MP and symmetric cache-only Shift.  These are mathematical interface
changes even when their literal implementation reuses most controls.

The coefficient-five and coefficient-six families change all formula
constructors rather than just the axiom order.  They include M5/M6,
N5/N6, H5/H6, Q5/Q6, QH5/QH6 and QN5/QN6.  A correct comparison must include
hereditary recoding of labels, unused residues, restricted-axiom bounds,
and the zero-code alias.  Several concrete radix specializations differ
by only one native decoder edge.  That is strong evidence of economical
implementation, but graph novelty alone does not establish a new
proof-search architecture.  The H5 report explicitly makes this
qualification about its own result.

\paragraph{Changing decoder association and physical returns.}
R and M4R reassociate the binary decoder while retaining the resolved
source computation.  The current corpus also contains 23 source-tree
rotations as research inputs.  Their presence is not evidence that each
has a separately established blank-start consistency theorem: source
binding, complete physical contracts and initialization are different
obligations.  J, F, A, X/CSZ, ZR, HA and HAZ exploit source registers known
to be zero to replace jumps by decrements, sometimes together with A/D
placement.  The required invariants may be small and elegant, but they
must hold at the actual physical entries, including initialization.

\paragraph{The newest interface change.}
AX11 pushes the final guarded head and falls through into the full MP
parameter reader.  It consequently consumes a fourth AX operand $z$
and produces an eleven-entry stack prefix.  Its completeness compiler
chooses $z=0$ where needed, while soundness must cover every $z$ because
the enumerator tries arbitrary candidates.  The new AX11H construction
retains that interface and changes its halt protocol.  The present
report's Pop analysis is needed precisely because these two inherited
mechanisms interact.  The contribution is a new composition with a new
relative compiler argument; it is not a newly invented numerical calculus.
\cite{ax11packet}

\subsection{Aliases, collisions of names and overlapping savings}
\label{cmp:aliases}

The same design often appears under several delivery names.  The following
identities are supported by the parsed tables, not inferred from similar
titles.  Full original paths and byte hashes remain in the local CSV/JSON
catalogues.\cite{repo,incoming}

\begin{table}[htbp]
\centering\small
\begin{tabularx}{\linewidth}{@{}p{0.29\linewidth}X@{}}
\toprule
Design or naming issue & Relationship in the supplied corpus\\
\midrule
Principal C & The same graph is delivered as the principal candidate in
five packets.\\
HS and lowercase NS & The cache-only NS in the lowercase packet equals HS.\\
N6 & Two principal filed deliveries reproduce the same design.\\
QN and NQ & Two names for the same principal graph.\\
QNS & The two filed QNS deliveries have the same principal table.\\
CSZ and X & The later X name denotes the existing CSZ graph.\\
QH5 & The AX11 companion and the separate incoming QH5 delivery have
identical dense bytes.\\
QN6 & The incoming QN6 delivery and the QN5/QN6 packet's companion
have identical dense bytes.\\
\bottomrule
\end{tabularx}
\caption{Repeated names do not count as repeated inventions.}
\label{cmp:alias-table}
\end{table}

The opposite error is equally serious: the sum-guard NS in
\path|PA194NS_research| is different from the cache-only NS/HS table.
The state number and an informal suffix therefore cannot identify a
machine.  Our tables retain both the graph identity and its historical
name, and record whether a start was declared or merely inferred from
the dense-numbering convention.

State savings also do not add automatically.  Existing SUMS and radix
folds both use native control 41, so their nominal benefits can overlap.
A grammar with fewer source nodes can require more physical readers or
invalidate a fold.  A local fold that works in isolation may fail when
another rewritten control occurs in its detour.  The correct unit of
verification is the final simultaneously modified table together with
its source binding and complete startup, not a sum of separate heuristic
scores.  This explains why some plausible combined optimizations remain
at 194 even though each ingredient was originally described as a saving.

\subsection{Why the stored 193-state objects are not solutions}
\label{cmp:sub194}

The current recognizer finds 53 distinct stored 193-state fingerprints.
That much larger number than the two principal earlier prototypes comes
from later runtime transplants and startup experiments.  None is supplied
as an established blank-start PA-consistency machine at the required
iPRA proof level.  Their separate evidence is retained in
\path|comparison/tables193_status.json|; the following are the decisive
families.\cite{incoming,ax11packet}

\paragraph{The coefficient-three shortcut is logically wrong.}
With the unchanged heads and the indicated coefficient-three encoding,
the collision $E_3(1,0)=I_3(0,1)=6$ permits a derivation of contradiction
code one.  The supplied expanded certificate has 73,393 primitive proof
steps and 116,171 shared arithmetic-expression nodes.  A later report
extends the obstruction to every nonnegative constant equality offset
within its precisely specified retained-head family.  This defeats that
small implementation; it is not a theorem excluding every alternative
three-residue encoding or every 193-state consistency machine.

\paragraph{The marker quotient has a runtime interface but no valid bootstrap.}
Identifying native control 301 with 88 and setting
$(88,0)\mapsto(0,L,51)$ and $(51,0)\mapsto(0,R,122)$ gives the widely
reused runtime193 core.  Its source contracts and changed marker identity
are separate from its blank initialization.  The original start reaches
an outward blank ray.  Later pair and triple repair screens report no
complete calibrated frame, but some runs exhaust finite search bounds;
those residual cases must not be described as proved divergent.

The current QN6 report contains a 49-table transplant suite.  Its recorded
all-start tests end in halts or certified rays, with no accepted frame.
A direct transfer to the new QN6 table reaches the same outward ray after
92 steps.  The ZFC-inspired relay experiment fails nine source outcomes
and has a 116-step blank-ray prefix.  The numbers in these reports are
attributed experimental evidence; this comparison does not silently
present every historical search as freshly rerun.

\paragraph{Restricted obstructions are useful without being global bounds.}
The direct-selector/disjoint-reader-fold study provides an exact dual
certificate limiting net savings to seven for its fixed 201-state parent
and specified rewrite family.  Overlapping folds, helpers targeting helpers,
new grammars and longer simulations lie outside that result.  Existing
work also includes radius-eight context screens, source-invariant
specialization, common-action mergers, incoming-edge bypasses, history
marking, all head orders for a fixed grammar, and extensive timed grammar
searches.  Negative outcomes constrain those searches; they do not prove
that 194 is minimal.  Repeating a weaker screen without a new invariant
would add little evidence.

\subsection{The weak-arithmetic requirement remains substantive}
\label{cmp:proof-status}

The requested endpoint is an equivalence provable in iPRA.  The current
PA README explicitly states that neither machine-specific direction has
yet been established there for primary PA194, original PA195 or PA196.
It records substantial completed intermediate results---including ordinary
finite-run simulation and modeled-scheduler components---while identifying
the remaining physical and proof-translation obligations.
The current incoming packets likewise distinguish ordinary written PA
calibrations from completed iPRA endpoints.\cite{pa-readme,incoming}

Two different distinctions must be preserved.  A mathematical proof can
be valid without having been compiled by Lean or Rocq.  But a proof that
uses PA does not become a proof in the weaker system iPRA merely because
its source code is small, its certificates are executable, or its witness
functions are described as primitive recursive.  The relevant correctness
statements and their internalization still require proof.  Conversely,
unfinished proof-assistant infrastructure is not itself a refutation of
the ordinary PA arguments in the historical papers.

This report therefore uses the corpus to establish exact construction
ancestry, structural distinction and explicit relative proof obligations.
The fresh AX11H verification supports its transition table, startup,
source interface and changed command compiler at the stated levels.
It does not manufacture the missing iPRA calibration by composing an
unproved premise with a verified local change.  The complete delivery and
frontier appendices make that limitation auditable for every supplied
design, rather than hiding it behind a single headline state count.
